\documentclass[12pt,a4paper]{article}

\usepackage{fontspec}
\setmainfont{Times New Roman}
\setmonofont{Courier New}

\usepackage[russian]{babel}
\usepackage[left=2.5cm,right=1.5cm,top=2cm,bottom=2cm]{geometry}
\usepackage{hyperref}
\usepackage{enumitem}
\usepackage{tabularx}
\usepackage{booktabs}
\usepackage{xcolor}
\usepackage{titlesec}
\usepackage{setspace}
\usepackage{amsmath}
\usepackage{amssymb}
\usepackage{microtype}
\usepackage{fancyhdr}
\usepackage{longtable}

\onehalfspacing

\hypersetup{
    colorlinks=true,
    linkcolor=black,
    urlcolor=blue,
    citecolor=black,
    pdftitle={Спецификация требований к ПО (SRS): Web-органайзер репетитора},
    pdfauthor={V01dC0Re}
}

\pagestyle{fancy}
\fancyhf{}
\fancyhead[L]{\small\textit{SRS: Web-органайзер и оркестратор учебного процесса репетитора}}
\fancyhead[R]{\small\textit{Версия 1.0}}
\fancyfoot[C]{\thepage}
\renewcommand{\headrulewidth}{0.4pt}
\renewcommand{\footrulewidth}{0pt}

\titleformat{\section}{\Large\bfseries}{\thesection.}{0.5em}{}
\titleformat{\subsection}{\large\bfseries}{\thesubsection.}{0.5em}{}
\titleformat{\subsubsection}{\normalsize\bfseries}{\thesubsubsection.}{0.5em}{}

\titlespacing*{\section}{0pt}{1.2em}{0.4em}
\titlespacing*{\subsection}{0pt}{0.8em}{0.3em}
\titlespacing*{\subsubsection}{0pt}{0.6em}{0.2em}

\setlist{nosep, topsep=2pt, itemsep=2pt, leftmargin=1.5em}

\newcolumntype{Y}{>{\raggedright\arraybackslash}X}
\newcolumntype{P}[1]{>{\raggedright\arraybackslash}p{#1}}

\begin{document}

\begin{titlepage}
    \centering
    \vspace*{1.5cm}
    
    {\Large\bfseries МИНИСТЕРСТВО НАУКИ И ВЫСШЕГО ОБРАЗОВАНИЯ РОССИЙСКОЙ ФЕДЕРАЦИИ\par}
    {\large УНИВЕРСИТЕТ ИТМО\par}
    
    \vspace{3.5cm}
    
    {\huge\bfseries СПЕЦИФИКАЦИЯ ТРЕБОВАНИЙ\par}
    \vspace{0.3cm}
    {\LARGE\bfseries К ПРОГРАММНОМУ ОБЕСПЕЧЕНИЮ\par}
    \vspace{0.5cm}
    {\Large (Software Requirements Specification — SRS)\par}
    
    \vspace{1.5cm}
    
    {\Large\textbf{Проект: «Web-органайзер и оркестратор учебного процесса репетитора»}\par}
    
    \vspace{1cm}
    {\large Стандарт: IEEE Std 830-1998 / ISO/IEC/IEEE 29148:2018\par}
    {\large Редакция: 1.0 (MVP Scope)\par}
    
    \vfill
    
    \begin{flushright}
        \textbf{Разработчик:} V01dC0Re\\
        \textbf{Направление:} Информационные системы
    \end{flushright}
    
    \vspace{1cm}
    {\large Санкт-Петербург, 2026\par}
\end{titlepage}

\newpage
\tableofcontents
\newpage

\section{Введение}

\subsection{Назначение документа}
Настоящий документ представляет собой спецификацию требований к программному обеспечению (Software Requirements Specification, SRS) веб-системы «Web-органайзер и оркестратор учебного процесса репетитора». Документ регламентирует функциональные и нефункциональные требования к системе, описывает архитектурные ограничения, внешние интерфейсы и модели взаимодействия. Предназначен для разработчиков, проектировщиков, тестировщиков и проверяющей комиссии.

\subsection{Область применения (Scope)}
Программный продукт представляет собой специализированную веб-систему (SaaS/On-Premise решение), ориентированную на частных преподавателей и репетиторов с высокой интенсивностью занятий. Основной целью системы является минимизация операционных и когнитивных издержек преподавателя посредством автоматизации:
\begin{itemize}
    \item гибкого диспетчирования временных слотов и расписания занятий;
    \item адаптивного управления учебными программами с каскадной синхронизацией;
    \item агрегации ежедневного цикла подготовки к урокам (Pre-Lesson Dashboard);
    \item централизованного хранения и дистрибуции учебно-методических материалов;
    \item консолидации коммуникационной и графической инфраструктуры урока в едином интерфейсе.
\end{itemize}

\subsection{Определения, акронимы и сокращения}
\begin{itemize}
    \item \textbf{SRS} (Software Requirements Specification) — спецификация требований к ПО.
    \item \textbf{MVP} (Minimum Viable Product) — минимально жизнеспособный продукт.
    \item \textbf{Curriculum Engine} — программный модуль адаптивного планирования и управления цепочками учебных тем.
    \item \textbf{Smart Slot Finder} — алгоритмический модуль поиска доступных временных окон с учетом буферных интервалов.
    \item \textbf{Буферное время} — технологический интервал времени между последовательными уроками, предназначенный для отдыха преподавателя и организационного переключения.
    \item \textbf{Спаренный урок} — учебное занятие удвоенной продолжительности, объединяющее две логические темы учебного плана.
    \item \textbf{Lesson Workspace} — контекстная панель занятия, предоставляющая мгновенный доступ к материалам, доске и видеосвязи.
    \item \textbf{Квотирование} — программное ограничение дискового пространства, выделяемого одному пользователю.
\end{itemize}

\subsection{Краткий обзор структуры}
Раздел 2 описывает общее назначение системы, портрет пользователей и эксплуатационные ограничения. Раздел 3 содержит формализованные функциональные требования к ключевым модулям (FR-1 -- FR-5). Раздел 4 регламентирует требования к интерфейсам. Раздел 5 определяет нефункциональные требования и показатели надежности. Раздел 6 содержит концептуальную модель данных, а Раздел 7 — матрицу трассируемости.

\section{Общее описание системы}

\subsection{Контекст продукта}
Система функционирует как независимое веб-приложение, доступное через современные браузеры. Архитектурно система объединяет в себе функции персонального календаря, специализированной LMS (Learning Management System) и агрегатора учебных инструментов (доски, видеоконференции, файлы), исключая необходимость ручной синхронизации между разрозненными сервисами.

\subsection{Классы и характеристики пользователей}
\begin{itemize}
    \item \textbf{Преподаватель (Репетитор)} — основной пользователь системы с правами полного управления: формирование расписания, создание мастер-шаблонов курсов, загрузка материалов, ведение чек-листов подготовки и организация рабочего пространства.
    \item \textbf{Системный администратор} — обслуживание серверной инфраструктуры, мониторинг дисковых квот, журналирование ошибок.
\end{itemize}

\subsection{Операционная среда и ограничения}
\begin{itemize}
    \item Клиентская часть: функционирование в актуальных версиях браузеров (Chromium $\ge$ 115, Firefox $\ge$ 118, Safari $\ge$ 16, Яндекс Браузер) на настольных и мобильных устройствах.
    \item Серверная часть: кроссплатформенная серверная среда, поддержка контейнеризации (Docker).
    \item Сетевые требования: обязательное шифрование трафика по протоколу HTTPS/TLS 1.3.
    \item Дисковая квота: ограничение дискового пространства в размере 1 Гб на учетную запись преподавателя в рамках базовой конфигурации.
\end{itemize}

\subsection{Предположения и зависимости}
\begin{itemize}
    \item Доступность сторонних сервисов видеоконференций (Zoom, Яндекс Телемост, Google Meet, VK Звонки) регулируется их собственными SLA и не гарантируется системой напрямую.
    \item Браузер пользователя поддерживает технологии HTML5 Canvas и WebSockets для работы встроенной графической доски.
\end{itemize}

\newpage

\section{Функциональные требования}

\subsection{Модуль 1. Интеллектуальное расписание и адаптивный календарь (SCHED)}

\subsubsection*{FR-1.1: Управление сеткой регулярных занятий}
\begin{itemize}
    \item \textbf{Описание:} Система должна обеспечивать создание циклических слотов расписания с еженедельным интервалом повторения.
    \item \textbf{Входные данные:} Идентификатор ученика, день недели, время начала, длительность (минуты), дата старта цикла, дата завершения цикла (опционально).
    \item \textbf{Бизнес-логика:} Генерация последовательности занятий в календаре на заданный горизонт планирования (по умолчанию 3 месяца вперед).
    \item \textbf{Выходные данные:} Сформированная сетка занятий в календарном представлении (неделя/месяц).
\end{itemize}

\subsubsection*{FR-1.2: Управление разовыми занятиями}
\begin{itemize}
    \item \textbf{Описание:} Возможность ручного назначения единичного занятия вне регулярного расписания.
    \item \textbf{Входные данные:} Идентификатор ученика, точная дата, время начала, длительность.
    \item \textbf{Выходные данные:} Карточка разового занятия в календаре со специальным визуальным маркером.
\end{itemize}

\subsubsection*{FR-1.3: Помощник подбора окон (Smart Slot Finder)}
\begin{itemize}
    \item \textbf{Описание:} Алгоритмический поиск и предложение оптимальных свободных интервалов времени для добавления нового ученика или переноса урока.
    \item \textbf{Входные данные:} Желаемая длительность урока $D$, диапазон дат/времени, настраиваемое буферное время $B$ (диапазон 0--60 мин, по умолчанию 15 мин).
    \item \textbf{Бизнес-логика:} Интервал $[t_{\text{start}}, t_{\text{end}}]$ считается валидным слотом, если:
    \begin{enumerate}
        \item $t_{\text{end}} - t_{\text{start}} = D$;
        \item Для любого существующего занятия $[s_i, e_i]$ выполняется: $t_{\text{start}} \ge e_i + B$ либо $t_{\text{end}} + B \le s_i$;
        \item Интервал укладывается в глобальные рабочие часы преподавателя.
    \end{enumerate}
    \item \textbf{Выходные данные:} Ранжированный список свободных окон с возможностью бронирования в один клик.
\end{itemize}

\subsubsection*{FR-1.4: Инфраструктурная карточка занятия}
\begin{itemize}
    \item \textbf{Описание:} Каждое занятие в календаре должно иметь интерактивную карточку быстрого запуска.
    \item \textbf{Элементы карточки:} Прямая кнопка-ссылка на сервис видеоконференции (Zoom / Яндекс Телемост / Google Meet / VK Звонки), кнопка перехода к рабочей доске, статус занятия, назначенная тема учебного плана.
\end{itemize}

\subsubsection*{FR-1.5: Статусная модель занятия}
\begin{itemize}
    \item \textbf{Описание:} Поддержка жизненного цикла занятия со следующими статусами: \texttt{Запланировано}, \texttt{Проведено}, \texttt{Отменено}, \texttt{Перенесено}.
\end{itemize}

\subsection{Модуль 2. Интерактивный учебный план (CURR)}

\subsubsection*{FR-2.1: Конструктор мастер-шаблонов курсов}
\begin{itemize}
    \item \textbf{Описание:} Создание и редактирование типовых образовательных программ с иерархической или линейной структурой.
    \item \textbf{Входные данные:} Название курса (например, «ОГЭ по математике: геометрия»), описание, перечень тем с указанием порядкового номера, описания и расчетного количества академических часов.
    \item \textbf{Выходные данные:} Сохраненный в библиотеке мастер-шаблон, доступный для тиражирования.
\end{itemize}

\subsubsection*{FR-2.2: Персонализация программы под ученика}
\begin{itemize}
    \item \textbf{Описание:} Создание индивидуальной копии мастер-шаблона для конкретного ученика.
    \item \textbf{Функционал:}
    \begin{itemize}[label=---]
        \item Изменение порядка тем посредством интерфейса drag-and-drop;
        \item Назначение приоритетов темам (высокий/базовый/факультативный);
        \item Объединение двух смежных тем в один спаренный урок;
        \item Добавление индивидуальных авторских тем преподавателя.
    \end{itemize}
\end{itemize}

\subsubsection*{FR-2.3: Каскадная реакция на отмену занятия}
\begin{itemize}
    \item \textbf{Описание:} Автоматическая адаптация расписания тем при срыве урока.
    \item \textbf{Бизнес-логика:} При смене статуса занятия на \texttt{Отменено}:
    \begin{enumerate}
        \item Тема, назначенная на отмененное занятие, автоматически открепляется от текущего слота;
        \item Данная тема назначается на следующее по хронологии запланированное занятие этого ученика;
        \item Все последующие темы цепочки сдвигаются на один плановый слот вперед (каскадный сдвиг).
    \end{enumerate}
    \item \textbf{Критерий приемки:} Целостность последовательности тем сохраняется, пропущенная тема не теряется из учебного плана.
\end{itemize}

\subsubsection*{FR-2.4: Адаптивная реакция на перенос занятия}
\begin{itemize}
    \item \textbf{Описание:} Обработка переноса занятия на нестандартное время.
    \item \textbf{Бизнес-логика:} Преподаватель выбирает новую дату/время (с опциональным использованием Smart Slot Finder). Текущая тема жестко закрепляется за новым временным слотом, при этом общая очередность остальных тем в основном регулярном графике не нарушается.
\end{itemize}

\subsection{Модуль 3. Панель подготовки к урокам (PREP)}

\subsubsection*{FR-3.1: Суточный дашборд готовности (Daily Prep View)}
\begin{itemize}
    \item \textbf{Описание:} Агрегирующий виджет, отображающий список уроков на ближайшие 24 часа (фильтрация: «Сегодня» / «Завтра»).
    \item \textbf{Представление:} Список карточек занятий, сгруппированных по хронологии, с интегрированным чек-листом готовности к каждому занятию.
\end{itemize}

\subsubsection*{FR-3.2: Формализованный 4-позиционный чек-лист}
\begin{itemize}
    \item \textbf{Описание:} Для каждого предстоящего урока система генерирует стандартизированный чек-лист готовности из 4 критериев:
    \begin{enumerate}
        \item \textbf{CL-1:} Проверено домашнее задание с предыдущего урока;
        \item \textbf{CL-2:} Сформированы теоретические и практические материалы к занятию;
        \item \textbf{CL-3:} Проверена актуальность ссылок на видеосвязь и доску;
        \item \textbf{CL-4:} Сформулировано и подготовлено домашнее задание на следующий урок.
    \end{enumerate}
    \item \textbf{Взаимодействие:} Фиксация состояния каждого пункта (чекбокс) с сохранением времени изменения.
\end{itemize}

\subsubsection*{FR-3.3: Индикатор совокупного прогресса дня}
\begin{itemize}
    \item \textbf{Описание:} Наглядное отображение степени завершенности подготовки к учебному дню.
    \item \textbf{Бизнес-логика:} Вычисление процента готовности по формуле:
    \[
    P = \left( \frac{\sum_{i=1}^{N} C_i}{4 \cdot N} \right) \times 100\%,
    \]
    где $N$ — количество уроков в выбранных сутках, $C_i \in [0, 4]$ — количество закрытых пунктов чек-листа $i$-го урока.
    \item \textbf{Визуализация:} Динамический Progress Bar и числовой индикатор.
\end{itemize}

\subsection{Модуль 4. Файловое хранилище и раздаточные материалы (FILE)}

\subsubsection*{FR-4.1: Мультиформатная загрузка материалов}
\begin{itemize}
    \item \textbf{Описание:} Загрузка пользовательских файлов с валидацией типа содержимого (MIME-type).
    \item \textbf{Поддерживаемые форматы:} Документы (\texttt{.pdf}, \texttt{.docx}, \texttt{.txt}), исходный код и скрипты (\texttt{.py}), растровые изображения (\texttt{.png}, \texttt{.jpg}, \texttt{.jpeg}).
    \item \textbf{Ограничение единичного файла:} Максимальный размер одного загружаемого файла — не более 50 Мб.
\end{itemize}

\subsubsection*{FR-4.2: Двухуровневая контекстная привязка файлов}
\begin{itemize}
    \item \textbf{Описание:} Возможность ассоциации файла с двумя уровнями абстракции:
    \begin{enumerate}
        \item \textbf{Привязка к теме учебного плана:} Файл является постоянным дидактическим материалом темы (доступен во всех уроках по этой теме).
        \item \textbf{Привязка к конкретному занятию:} Файл является оперативным материалом (результаты решения, скан работы ученика на данном уроке).
    \end{enumerate}
\end{itemize}

\subsubsection*{FR-4.3: Планшет урока (Lesson Workspace)}
\begin{itemize}
    \item \textbf{Описание:} Автоматически активируемый блок «Материалы урока» в карточке занятия при наступлении даты урока.
    \item \textbf{Функционал:} Отображение всех прикрепленных к занятию и теме документов с возможностью их скачивания или предпросмотра в один клик без перехода в файловый менеджер.
\end{itemize}

\subsubsection*{FR-4.4: Квотирование дискового пространства и лимиты}
\begin{itemize}
    \item \textbf{Описание:} Контроль дискового пространства учетной записи преподавателя.
    \item \textbf{Бизнес-логика:}
    \begin{itemize}[label=---]
        \item Базовый лимит дискового пространства: ровно \textbf{1 Гб} на пользователя;
        \item Постоянный пересчет суммарного объема файлов $\sum \text{FileSize}$;
        \item Визуализация в интерфейсе в виде прогресс-бара и счетчика «Занято $X$ Мб из 1024 Мб»;
        \item При попытке загрузки файла, превышающего остаток квоты: блокировка операции, вывод предупреждающего модального окна с сообщением об исчерпании лимита и демонстрационным предложением расширения тарифа (демо-подписка).
    \end{itemize}
\end{itemize}

\subsection{Модуль 5. Интеграция рабочего пространства (доски и связь) (WRK)}

\subsubsection*{FR-5.1: Привязка внешних интерактивных досок}
\begin{itemize}
    \item \textbf{Описание:} Интеграция постоянных внешних онлайн-досок в профиль ученика.
    \item \textbf{Поддерживаемые платформы:} Miro, sBoard, Conceptboard, произвольные URL.
    \item \textbf{Функционал:} Сохранение URL-адреса персональной доски ученика; автоматическая подстановка активной ссылки во все карточки занятий данного ученика.
\end{itemize}

\subsubsection*{FR-5.2: Встроенная легковесная доска (Embedded Canvas)}
\begin{itemize}
    \item \textbf{Описание:} Встроенный компонент быстрой графической доски на базе открытых библиотек (Excalidraw / tldraw).
    \item \textbf{Назначение:} Оперативные математические расчеты, запись формул, построение быстрых схем без необходимости регистрации преподавателя и ученика на сторонних сервисах.
    \item \textbf{Функционал:}
    \begin{itemize}[label=---]
        \item Рисование от руки, геометрические примитивы, текст, стрелки;
        \item Быстрый экспорт созданного фрагмента в изображение (\texttt{.png}) с автоматическим прикреплением к текущему уроку в Lesson Workspace.
    \end{itemize}
\end{itemize}

\newpage

\section{Требования к внешним интерфейсам}

\subsection{Пользовательский интерфейс (UI/UX)}
\begin{itemize}
    \item \textbf{Адаптивность:} Интерфейс должен корректно масштабироваться под разрешения экранов от $360\times640$ (смартфоны) до $1920\times1080$ и выше (десктоп).
    \item \textbf{Цветовое кодирование:} Статусы занятий, индикаторы готовности чек-листов и прогресс-бар квоты диска должны иметь однозначную цветовую дифференциацию (зеленый — готово/успех, желтый — в процессе/внимание, красный — срыв/превышение).
    \item \textbf{Эргономика:} Доступ к ключевым операциям (запуск звонка, открытие материалов, отметка чек-листа) должен осуществляться не более чем в 2 клика с главной страницы.
\end{itemize}

\subsection{Программные интерфейсы (API)}
\begin{itemize}
    \item Внутренний обмен между клиентской и серверной частью должен строиться по архитектурному стилю REST API с сериализацией данных в формате JSON.
    \item Для обновления статусов уроков и работы встроенной доски в реальном времени рекомендуется использование протокола WebSocket.
\end{itemize}

\subsection{Интерфейсы сторонних сервисов}
\begin{itemize}
    \item Прямая генерация ссылок протоколов (deep linking) для быстрого запуска настольных приложений видеосвязи (\texttt{zoommtg://}, \texttt{telemost://} и веб-ссылок Google Meet).
\end{itemize}

\section{Нефункциональные требования}

\subsection{Требования к производительности}
\begin{itemize}
    \item Время отклика сервера на базовые CRUD-запросы (создание слота, смена статуса) не должно превышать 300 мс при 95-м процентиле (p95).
    \item Время пересчета расписания при каскадном сдвиге цепочки тем (до 100 тем) не должно превышать 500 мс.
    \item Время первоначальной загрузки страницы календаря при скорости сети $\ge$ 10 Мбит/с — не более 1.5 с.
\end{itemize}

\subsection{Надежность и доступность}
\begin{itemize}
    \item Коэффициент доступности системы (Availability): не менее 99.5\% в рабочее время (с 08:00 до 23:00 МСК).
    \item Транзакционность: операции сдвига учебного плана и обновления слотов должны выполняться в рамках атомарных транзакций базы данных для предотвращения рассинхронизации.
    \item Резервное копирование: автоматическое создание снапшотов БД не реже одного раза в 24 часа.
\end{itemize}

\subsection{Требования к безопасности и защите данных}
\begin{itemize}
    \item Аутентификация преподавателей с использованием безопасных хэшей паролей (Argon2id или bcrypt) и сессионных JWT-токенов.
    \item Разграничение прав доступа: строгая изоляция данных (multi-tenant логика на уровне идентификатора преподавателя \texttt{tutor\_id}). Преподаватель не имеет технической возможности получить доступ к расписанию или файлам другого преподавателя.
    \item Защита от типовых веб-уязвимостей: валидация входных данных, экранирование SQL-инъекций (через ORM/параметризованные запросы), защита от CSRF и XSS атак.
\end{itemize}

\newpage

\section{Концептуальная модель данных}

Ниже приведено формальное описание ключевых сущностей базы данных системы:

\begin{table}[h!]
\centering
\small
\begin{tabularx}{\textwidth}{l P{3.5cm} Y}
\toprule
\textbf{Сущность} & \textbf{Назначение} & \textbf{Ключевые атрибуты} \\
\midrule
\texttt{Tutor} & Учетная запись преподавателя & \texttt{id}, \texttt{email}, \texttt{password\_hash}, \texttt{storage\_used}, \texttt{created\_at} \\
\addlinespace
\texttt{Student} & Профиль ученика & \texttt{id}, \texttt{tutor\_id}, \texttt{full\_name}, \texttt{board\_url}, \texttt{default\_meeting\_url} \\
\addlinespace
\texttt{CourseTemplate} & Мастер-шаблон учебного курса & \texttt{id}, \texttt{tutor\_id}, \texttt{title}, \texttt{description}, \texttt{is\_public} \\
\addlinespace
\texttt{Topic} & Тема курса / индивидуального плана & \texttt{id}, \texttt{template\_id}, \texttt{student\_id}, \texttt{order\_index}, \texttt{title}, \texttt{hours\_count} \\
\addlinespace
\texttt{Lesson} & Конкретное занятие в расписании & \texttt{id}, \texttt{tutor\_id}, \texttt{student\_id}, \texttt{start\_time}, \texttt{end\_time}, \texttt{status}, \texttt{topic\_id} \\
\addlinespace
\texttt{PreLessonCheck} & Чек-лист готовности к занятию & \texttt{id}, \texttt{lesson\_id}, \texttt{hw\_checked}, \texttt{materials\_ready}, \texttt{links\_verified}, \texttt{next\_hw\_set} \\
\addlinespace
\texttt{MaterialFile} & Загруженный дидактический файл & \texttt{id}, \texttt{tutor\_id}, \texttt{file\_name}, \texttt{file\_size}, \texttt{mime\_type}, \texttt{topic\_id}, \texttt{lesson\_id} \\
\bottomrule
\end{tabularx}
\caption{Ключевые сущности концептуальной модели данных}
\end{table}

\section{Матрица трассируемости требований}

\begin{table}[h!]
\centering
\small
\begin{tabularx}{\textwidth}{l Y l}
\toprule
\textbf{Код требования} & \textbf{Наименование требования} & \textbf{Связанные сущности БД} \\
\midrule
FR-1.1, FR-1.2 & Сетка регулярных и разовых занятий & \texttt{Lesson}, \texttt{Student} \\
FR-1.3 & Smart Slot Finder & \texttt{Lesson}, \texttt{Tutor} \\
FR-1.4, FR-1.5 & Карточка урока и статусы & \texttt{Lesson}, \texttt{Student} \\
FR-2.1 & Мастер-шаблоны курсов & \texttt{CourseTemplate}, \texttt{Topic} \\
FR-2.2 & Персонализация программы & \texttt{Topic}, \texttt{Student} \\
FR-2.3, FR-2.4 & Каскадный сдвиг при отмене/переносе & \texttt{Lesson}, \texttt{Topic} \\
FR-3.1, FR-3.2, FR-3.3 & Дашборд подготовки и чек-листы & \texttt{Lesson}, \texttt{PreLessonCheck} \\
FR-4.1, FR-4.2, FR-4.3 & Загрузка и контекстная привязка файлов & \texttt{MaterialFile}, \texttt{Topic}, \texttt{Lesson} \\
FR-4.4 & Квотирование хранилища (1 Гб) & \texttt{Tutor}, \texttt{MaterialFile} \\
FR-5.1, FR-5.2 & Интеграция досок и встраиваемый холст & \texttt{Student}, \texttt{Lesson}, \texttt{MaterialFile} \\
\bottomrule
\end{tabularx}
\caption{Матрица трассируемости функциональных требований}
\end{table}

\end{document}
